\section{Theory}
\markright{Theory}

Computer-aided design (CAD) software has become an indispensable tool in modern mechanical engineering for creating, analyzing, and documenting mechanical components and assemblies \cite{lab03_solidworks_guide}. SolidWorks, a parametric 3D modeling platform, enables engineers to design components using feature-based modeling where geometric constraints and relationships control part dimensions. This parametric approach allows for rapid design iterations and ensures dimensional consistency across related components. The fundamental workflow involves sketching 2D profiles on selected planes, applying geometric and dimensional constraints to create the sketch, and then utilizing solid modeling features (such as extrude, revolve, sweep, and loft) to generate 3D part geometry. Each feature maintains a history-based relationship, allowing modifications to upstream features to automatically propagate through the design, thereby facilitating efficient design refinement and variant generation.

Assembly design in SolidWorks involves combining individual parts into a functional assembly through the application of kinematic constraints that define part relationships and restrict degrees of freedom \cite{lab03_mech_design}. Standard assembly constraints include coincident, parallel, perpendicular, and distance constraints, which together simulate the physical constraints and motion of real mechanical systems. Once parts are assembled with appropriate constraints, kinematic analysis can verify that components move as intended without interference or constraint conflicts. Technical documentation, including detail drawings with dimensions, tolerances, and notes, along with assembly drawings and bills of materials, constitutes the complete design deliverable \cite{lab03_cad_design}. This documentation serves as the communication bridge between design intent and manufacturing execution, ensuring that fabricated components meet functional and quality requirements.
